% =============================================================================
% TryLM Tutorial Book - Preamble
% =============================================================================

% -----------------------------------------------------------------------------
% Document Class Settings
% -----------------------------------------------------------------------------
% Using KOMA-Script book class for better typography
\usepackage{scrhack}

% -----------------------------------------------------------------------------
% Page Layout
% -----------------------------------------------------------------------------
\usepackage[margin=1in]{geometry}
\usepackage{fancyhdr}

\pagestyle{fancy}
\fancyhf{}
\fancyhead[LE,RO]{\thepage}
\fancyhead[LO]{\nouppercase{\rightmark}}
\fancyhead[RE]{\nouppercase{\leftmark}}
\renewcommand{\headrulewidth}{0.4pt}

% -----------------------------------------------------------------------------
% Typography and Fonts
% -----------------------------------------------------------------------------
\usepackage[T1]{fontenc}
\usepackage{lmodern}
\usepackage{microtype}

% -----------------------------------------------------------------------------
% Mathematics
% -----------------------------------------------------------------------------
\usepackage{amsmath}
\usepackage{amssymb}
\usepackage{amsthm}
\usepackage{mathtools}
\usepackage{bm}  % Bold math symbols

% Custom math commands
\newcommand{\softmax}{\operatorname{softmax}}
\newcommand{\ReLU}{\operatorname{ReLU}}
\newcommand{\sigmoid}{\operatorname{sigmoid}}
% \tanh is already defined in LaTeX
\newcommand{\LayerNorm}{\operatorname{LayerNorm}}
\newcommand{\RMSNorm}{\operatorname{RMSNorm}}
\newcommand{\Attention}{\operatorname{Attention}}
\newcommand{\FFN}{\operatorname{FFN}}
\newcommand{\SwiGLU}{\operatorname{SwiGLU}}
\newcommand{\Swish}{\operatorname{Swish}}
\newcommand{\transpose}{\mathsf{T}}
\newcommand{\vocab}{\mathcal{V}}
\newcommand{\loss}{\mathcal{L}}

% -----------------------------------------------------------------------------
% Code Listings
% -----------------------------------------------------------------------------
\usepackage{listings}
\usepackage{xcolor}

% Define colors for code highlighting
\definecolor{codebg}{RGB}{248, 248, 248}
\definecolor{codeframe}{RGB}{200, 200, 200}
\definecolor{codecomment}{RGB}{106, 153, 85}
\definecolor{codekeyword}{RGB}{86, 156, 214}
\definecolor{codestring}{RGB}{206, 145, 120}
\definecolor{codenumber}{RGB}{181, 206, 168}
\definecolor{codefunction}{RGB}{220, 220, 170}

% Python style
\lstdefinestyle{python}{
    language=Python,
    basicstyle=\small\ttfamily,
    backgroundcolor=\color{codebg},
    frame=single,
    framesep=5pt,
    rulecolor=\color{codeframe},
    commentstyle=\color{codecomment}\itshape,
    keywordstyle=\color{codekeyword}\bfseries,
    stringstyle=\color{codestring},
    numberstyle=\tiny\color{codenumber},
    numbers=left,
    numbersep=8pt,
    breaklines=true,
    breakatwhitespace=true,
    showstringspaces=false,
    tabsize=4,
    captionpos=b,
    morekeywords={self, True, False, None, yield, async, await},
    emph={__init__, forward, __call__, __len__, __getitem__},
    emphstyle=\color{codefunction},
    literate={->}{$\rightarrow$}1,
}

% Shell/Bash style
\lstdefinestyle{shell}{
    language=bash,
    basicstyle=\small\ttfamily,
    backgroundcolor=\color{codebg},
    frame=single,
    framesep=5pt,
    rulecolor=\color{codeframe},
    commentstyle=\color{codecomment}\itshape,
    keywordstyle=\color{codekeyword}\bfseries,
    numbers=none,
    breaklines=true,
    showstringspaces=false,
    tabsize=4,
}

% YAML style
\lstdefinestyle{yaml}{
    basicstyle=\small\ttfamily,
    backgroundcolor=\color{codebg},
    frame=single,
    framesep=5pt,
    rulecolor=\color{codeframe},
    commentstyle=\color{codecomment}\itshape,
    keywordstyle=\color{codekeyword},
    stringstyle=\color{codestring},
    numbers=none,
    breaklines=true,
    showstringspaces=false,
    tabsize=2,
}

\lstset{style=python}

% Inline code command
\newcommand{\code}[1]{\texttt{\small #1}}
\newcommand{\pycode}[1]{\lstinline[style=python]{#1}}

% -----------------------------------------------------------------------------
% Diagrams with TikZ
% -----------------------------------------------------------------------------
\usepackage{tikz}
\usepackage{pgfplots}
\pgfplotsset{compat=1.18}
\usetikzlibrary{
    positioning,
    arrows.meta,
    shapes.geometric,
    shapes.misc,
    shapes.symbols,
    shapes.multipart,
    matrix,
    calc,
    fit,
    backgrounds,
    decorations.pathreplacing,
    decorations.markings,
    patterns,
    chains,
}

% TikZ styles for neural network diagrams
\tikzset{
    % Basic node styles
    neuron/.style={circle, draw, minimum size=8mm, inner sep=0pt},
    hidden neuron/.style={neuron, fill=blue!20},
    input neuron/.style={neuron, fill=green!20},
    output neuron/.style={neuron, fill=red!20},
    % Block styles
    block/.style={rectangle, draw, rounded corners=3pt, minimum height=10mm,
                  minimum width=20mm, align=center, fill=blue!10},
    attention block/.style={block, fill=orange!20},
    ffn block/.style={block, fill=green!20},
    norm block/.style={block, fill=yellow!20, minimum width=15mm},
    embed block/.style={block, fill=purple!20},
    % Arrow styles
    arrow/.style={-{Stealth[length=2mm]}, thick},
    skip arrow/.style={arrow, dashed, gray},
    data flow/.style={arrow, blue!70},
    % Matrix styles
    tensor/.style={matrix of nodes, nodes={draw, minimum size=6mm, anchor=center},
                   column sep=-\pgflinewidth, row sep=-\pgflinewidth},
    % Annotation styles
    annot/.style={font=\footnotesize, text=gray},
    label box/.style={rectangle, draw, dashed, gray, inner sep=2mm},
}

% -----------------------------------------------------------------------------
% Colored Boxes (tcolorbox)
% -----------------------------------------------------------------------------
\usepackage[most]{tcolorbox}

% Learning objectives box
\newtcolorbox{objectives}{
    enhanced,
    colback=blue!5,
    colframe=blue!75!black,
    fonttitle=\bfseries,
    title=Learning Objectives,
    before upper={\begin{itemize}[leftmargin=*,nosep]},
    after upper={\end{itemize}},
    boxrule=1pt,
    arc=3pt,
    left=5pt,
    right=5pt,
    top=5pt,
    bottom=5pt,
}

% Key concept box
\newtcolorbox{keypoint}[1][]{
    enhanced,
    colback=green!5,
    colframe=green!50!black,
    fonttitle=\bfseries,
    title={Key Concept},
    boxrule=1pt,
    arc=3pt,
    left=5pt,
    right=5pt,
    #1
}

% Common pitfall/warning box
\newtcolorbox{pitfall}[1][]{
    enhanced,
    colback=red!5,
    colframe=red!75!black,
    fonttitle=\bfseries,
    title={Common Pitfall},
    boxrule=1pt,
    arc=3pt,
    left=5pt,
    right=5pt,
    before upper={\setlength{\parindent}{0pt}},
    #1
}

% Exercise box
\newtcolorbox{exercise}[1][]{
    enhanced,
    colback=orange!5,
    colframe=orange!75!black,
    fonttitle=\bfseries,
    title={Exercise},
    boxrule=1pt,
    arc=3pt,
    left=5pt,
    right=5pt,
    #1
}

% Note/tip box
\newtcolorbox{note}[1][]{
    enhanced,
    colback=cyan!5,
    colframe=cyan!75!black,
    fonttitle=\bfseries,
    title={Note},
    boxrule=1pt,
    arc=3pt,
    left=5pt,
    right=5pt,
    #1
}

% Definition box
\newtcolorbox{definition}[1][]{
    enhanced,
    colback=gray!5,
    colframe=gray!75!black,
    fonttitle=\bfseries,
    title={Definition},
    boxrule=1pt,
    arc=3pt,
    left=5pt,
    right=5pt,
    #1
}

% Summary box (for end of chapter)
\newtcolorbox{summary}{
    enhanced,
    colback=violet!5,
    colframe=violet!75!black,
    fonttitle=\bfseries,
    title={Chapter Summary},
    boxrule=1pt,
    arc=3pt,
    left=5pt,
    right=5pt,
    before upper={\begin{itemize}[leftmargin=*,nosep]},
    after upper={\end{itemize}},
}

% -----------------------------------------------------------------------------
% Tables
% -----------------------------------------------------------------------------
\usepackage{booktabs}
\usepackage{array}
\usepackage{tabularx}
\usepackage{longtable}
\usepackage{multirow}

% Better column types
\newcolumntype{L}[1]{>{\raggedright\arraybackslash}p{#1}}
\newcolumntype{C}[1]{>{\centering\arraybackslash}p{#1}}
\newcolumntype{R}[1]{>{\raggedleft\arraybackslash}p{#1}}

% -----------------------------------------------------------------------------
% Figures and Captions
% -----------------------------------------------------------------------------
\usepackage{graphicx}
\usepackage{caption}
\usepackage{subcaption}
\usepackage{float}

\captionsetup{
    font=small,
    labelfont=bf,
    format=hang,
    margin=10pt,
}

% -----------------------------------------------------------------------------
% Cross-references and Links
% -----------------------------------------------------------------------------
\usepackage{hyperref}
\hypersetup{
    colorlinks=true,
    linkcolor=blue!70!black,
    citecolor=green!50!black,
    urlcolor=blue!70!black,
    bookmarks=true,
    bookmarksnumbered=true,
    pdfauthor={TryLM Tutorial},
    pdftitle={Building Language Models from Scratch},
    pdfsubject={Language Model Tutorial},
    pdfkeywords={LLM, Transformer, GPT, PyTorch, Deep Learning},
}

\usepackage{cleveref}
\crefname{figure}{Figure}{Figures}
\crefname{table}{Table}{Tables}
\crefname{equation}{Equation}{Equations}
\crefname{chapter}{Chapter}{Chapters}
\crefname{section}{Section}{Sections}
\crefname{appendix}{Appendix}{Appendices}

% -----------------------------------------------------------------------------
% Lists
% -----------------------------------------------------------------------------
\usepackage{enumitem}
\setlist{noitemsep, topsep=3pt}

% -----------------------------------------------------------------------------
% Bibliography (if needed)
% -----------------------------------------------------------------------------
\usepackage[style=numeric, sorting=none, backend=biber]{biblatex}
% \addbibresource{references.bib}

% -----------------------------------------------------------------------------
% Miscellaneous
% -----------------------------------------------------------------------------
\usepackage{xspace}
\usepackage{siunitx}  % SI units
\usepackage{datetime2}  % Date formatting

% Custom commands
\newcommand{\TryLM}{\textsc{TryLM}\xspace}
\newcommand{\PyTorch}{\textsc{PyTorch}\xspace}
\newcommand{\Python}{\textsc{Python}\xspace}
\newcommand{\GPT}{\textsc{GPT}\xspace}
\newcommand{\transformer}{Transformer\xspace}

% Parameter notation
\newcommand{\param}[1]{\texttt{#1}}
\newcommand{\file}[1]{\texttt{#1}}
\newcommand{\cmd}[1]{\texttt{\$ #1}}

% Dimension notation (e.g., [batch, seq, hidden])
\newcommand{\shape}[1]{{\small\textsf{[#1]}}}

% -----------------------------------------------------------------------------
% Chapter style customization
% -----------------------------------------------------------------------------
\usepackage{titlesec}

\titleformat{\chapter}[display]
    {\normalfont\huge\bfseries}
    {\chaptertitlename\ \thechapter}
    {20pt}
    {\Huge}

\titleformat{\section}
    {\normalfont\Large\bfseries}
    {\thesection}
    {1em}
    {}

\titleformat{\subsection}
    {\normalfont\large\bfseries}
    {\thesubsection}
    {1em}
    {}

% -----------------------------------------------------------------------------
% Algorithms (optional, for pseudocode)
% -----------------------------------------------------------------------------
\usepackage{algorithm}
\usepackage{algpseudocode}

\algrenewcommand\algorithmicrequire{\textbf{Input:}}
\algrenewcommand\algorithmicensure{\textbf{Output:}}

% -----------------------------------------------------------------------------
% End of Preamble
% -----------------------------------------------------------------------------
